%%% SDGFT-2026-01
%%% Axiomatic Foundations of Scale-Dependent Geometric Field Theory
%%%
\documentclass[amsmath,amssymb,aps,prd,twocolumn,superscriptaddress,nofootinbib]{revtex4-2}
\usepackage{graphicx}
\usepackage{hyperref}
\usepackage{xcolor}
\usepackage{booktabs}
\usepackage{float}
\usepackage{sdgft-macros}

\begin{document}

\preprint{SDGFT-2026-01}

\title{Axiomatic Foundations of Scale-Dependent Geometric Field Theory}

\author{David~A.~Besemer}
\email{david.besemer@sdgft.org}
\affiliation{Independent Researcher}
\date{\today}

\begin{abstract}
We present the axiomatic foundations of Scale-Dependent Geometric Field Theory
(SDGFT), a zero-parameter framework that derives the constants of Nature from
the geometry of the 24-cell polytope $\{3,4,3\}$.  Two axioms suffice: (I)~the
Fibonacci-lattice conflict parameter $\Delta = F_5/24 = 5/24$ and (II)~the
elementary lattice tension $\delta = 1/24$.  From these, the golden ratio
$\varphi = (1+\sqrt{5})/2$, the six-cone partition of $S^3$, and the critical
spacetime dimension $D^* = 67/24$ emerge without free parameters.  We establish
consistency checks, connect the construction to the Dedekind eta function and
the bosonic string dimension $D_{\text{bos}} = 26$, and contrast SDGFT's
zero-parameter philosophy with the landscape problem of string theory and the
Barbero--Immirzi ambiguity of loop quantum gravity.
\end{abstract}

\maketitle

%======================================================================
\section{Introduction}\label{sec:intro}
%======================================================================

The quest for a unified description of quantum gravity and the Standard Model
(SM) of particle physics has produced two major research programmes---string
theory and loop quantum gravity (LQG)---each carrying adjustable parameters
or unconstrained moduli.  String theory admits of order $10^{500}$ flux vacua
in its landscape~\cite{Bousso2000,Susskind2003}, while LQG contains the
Barbero--Immirzi parameter $\gamma_{\text{BI}}$ whose value is fixed only by
matching the Bekenstein--Hawking entropy~\cite{Immirzi1997,Rovelli1996}.

Scale-Dependent Geometric Field Theory (SDGFT) takes a radically different
approach: \emph{every} physical constant is derived from the combinatorial and
topological properties of a single geometric object, the 24-cell polytope
$\{3,4,3\}$.  The theory contains \emph{zero} adjustable parameters.  In this
paper---the first of the SDGFT Band~I series---we lay the axiomatic
foundations.

The programme proceeds as follows.  Section~\ref{sec:24cell} reviews the
properties of the 24-cell.  Sections~\ref{sec:axiom1} and~\ref{sec:axiom2}
state the two axioms.  Section~\ref{sec:emergent} derives the emergent
constants.  Section~\ref{sec:checks} presents consistency checks.
Section~\ref{sec:dedekind} connects the axioms to modular forms and the bosonic
string.  Section~\ref{sec:zero} contrasts SDGFT with existing approaches, and
Sec.~\ref{sec:conclusions} concludes.

%======================================================================
\section{The 24-Cell Polytope \texorpdfstring{$\{3,4,3\}$}{}}\label{sec:24cell}
%======================================================================

The 24-cell is the unique self-dual regular polytope in four dimensions that has
no three-dimensional analogue.  Its Schl\"afli symbol $\{3,4,3\}$ encodes the
local geometry: each cell is an octahedron $\{3,4\}$, four of which meet at
every edge.  The combinatorial data are summarised in Table~\ref{tab:24cell}.

\begin{table}[H]
\centering
\caption{Combinatorial data of the 24-cell.}\label{tab:24cell}
\begin{tabular}{@{}lc@{}}
\toprule
Element & Count \\
\midrule
Vertices ($f_0$) & 24 \\
Edges ($f_1$) & 96 \\
Triangular 2-faces ($f_2$) & 96 \\
Octahedral 3-cells ($f_3$) & 24 \\
\bottomrule
\end{tabular}
\end{table}

The Euler characteristic satisfies
\begin{equation}\label{eq:euler}
  \chi = f_0 - f_1 + f_2 - f_3 = 24 - 96 + 96 - 24 = 0\,,
\end{equation}
consistent with a closed 4-manifold tiling.  The 24 vertices may be placed on
the unit 3-sphere $S^3 \subset \mathbb{R}^4$ at the positions
\begin{equation}\label{eq:verts}
  \bigl\{\,(\pm1,0,0,0),\;
  \tfrac{1}{\sqrt{2}}(\pm1,\pm1,0,0)\,\bigr\}
\end{equation}
and all permutations thereof, giving 8 vertices of the form $(\pm1,0,0,0)$ and
16 of the form $\tfrac{1}{\sqrt{2}}(\pm1,\pm1,0,0)$.  The symmetry group is
the Weyl group of $F_4$, of order $|W(F_4)| = 1152$; see
Ref.~\cite{Besemer2026_02} for a detailed treatment.

A key structural property is the decomposition of the 24-cell into three
mutually orthogonal 16-cells (cross-polytopes), each contributing 8~vertices:
\begin{equation}\label{eq:decomp}
  24\text{-cell} = 16\text{-cell}_1 \cup 16\text{-cell}_2
                  \cup 16\text{-cell}_3\,,
  \quad 8+8+8=24\,.
\end{equation}
This triple decomposition is intimately connected with $D_4$ triality and
underlies the emergence of three fermion generations in SDGFT.

%======================================================================
\section{Axiom I: Fibonacci-Lattice Conflict}\label{sec:axiom1}
%======================================================================

\noindent\textbf{Axiom~I.}\quad\emph{The Fibonacci-lattice conflict parameter is}
\begin{equation}\label{eq:Delta}
  \boxed{\;\Delta \;=\; \frac{F_5}{24} \;=\; \frac{5}{24}\;\approx\; 0.20833\,.\;}
\end{equation}

Here $F_5 = 5$ is the fifth Fibonacci number.  The numerator encodes the
fivefold symmetry of the icosahedron---the symmetry incompatible with a
crystallographic lattice in three dimensions---while the denominator~$24$
counts the vertices (equivalently, cells) of the 24-cell.  The parameter
$\Delta$ therefore quantifies the geometric frustration arising from embedding
Fibonacci (icosahedral) order into a lattice defined by $\{3,4,3\}$.

Physically, $\Delta$ sets the fundamental mismatch between the continuous
rotational symmetry of the golden ratio and the discrete symmetry of the
crystallographic lattice.  In the language of condensed-matter physics, it is
analogous to a phason strain in a quasicrystal.

%======================================================================
\section{Axiom II: Elementary Lattice Tension}\label{sec:axiom2}
%======================================================================

\noindent\textbf{Axiom~II.}\quad\emph{The elementary lattice tension is}
\begin{equation}\label{eq:delta}
  \boxed{\;\delta \;=\; \frac{1}{24}\;\approx\; 0.04167\,.\;}
\end{equation}

This is the smallest non-zero eigenvalue of the lattice-strain operator on the
24-cell.  Because the 24-cell has exactly 24~cells, the unit of geometric
tension is $1/24$.  The parameter $\delta$ plays the role of a ``quantum'' of
curvature defect on $S^3$: the minimal angular deficit that a single vertex
insertion can produce.

Together, Axioms~I and~II fix every constant of the theory.  No further input
is required.

%======================================================================
\section{Emergent Constants}\label{sec:emergent}
%======================================================================

\subsection{The golden ratio}

The ratio of the two axioms yields
\begin{equation}\label{eq:ratio}
  \frac{\Delta}{\delta} = \frac{5/24}{1/24} = 5\,.
\end{equation}
The integer~5 uniquely singles out the golden ratio $\varphi = (1+\sqrt{5})/2$
among all algebraic numbers: it is the positive root of $x^2 - x - 1 = 0$
whose discriminant is~$\sqrt{5}$.  More precisely,
\begin{equation}\label{eq:phi}
  \varphi = \frac{1 + \sqrt{\Delta/\delta}}{2}
           = \frac{1 + \sqrt{5}}{2} \approx 1.61803\,.
\end{equation}

\subsection{Cone half-angle}

The sum of the two axioms yields
\begin{equation}\label{eq:sum}
  \Delta + \delta = \frac{5}{24} + \frac{1}{24} = \frac{6}{24}
                  = \frac{1}{4} = \sin^2 30°\,.
\end{equation}
The identification $\sin^2\theta_{\max} = 1/4$ fixes the maximal cone
half-angle $\theta_{\max} = 30°$.

\subsection{Six-cone partition}

A cone of half-angle $30°$ on $S^3$ subtends a solid angle
\begin{equation}\label{eq:solidangle}
  \Omega(\theta_{\max}) = 2\pi\bigl(1 - \cos\theta_{\max}\bigr)
  = 2\pi\!\left(1 - \frac{\sqrt{3}}{2}\right).
\end{equation}
The total solid angle of $S^3$ is $2\pi^2$.  The partition number is
\begin{equation}\label{eq:partition}
  N_{\text{cones}} = \frac{2\pi^2}{\Omega(\theta_{\max})}
  = \frac{2\pi^2}{2\pi(1 - \sqrt{3}/2)} = \frac{\pi}{1 - \sqrt{3}/2}\,.
\end{equation}
For the flat-space (small-angle) limit, $360°/60° = 6$, giving exactly six
congruent cones.  This six-cone partition is the geometric origin of the three
fermion generations $\times$ two chiralities; see Ref.~\cite{Besemer2026_03}
for the detailed proof.

\subsection{Critical dimension}

The critical spacetime dimension of SDGFT is
\begin{equation}\label{eq:Dstar}
  D^* = \frac{67}{24} \approx 2.79167 \quad(\text{tree level})\,,
\end{equation}
with the fixed-point value $D^* \approx 2.79676$ obtained by resumming
logarithmic corrections.  The spectral dimension at the Planck scale is
\begin{equation}\label{eq:spectral}
  n = \frac{D^*}{2} = \frac{67}{48} \approx 1.39583\,.
\end{equation}

%======================================================================
\section{Consistency Checks}\label{sec:checks}
%======================================================================

The two axioms satisfy a pair of exact arithmetic identities:
\begin{align}
  \Delta + \delta &= \frac{5}{24} + \frac{1}{24}
                   = \frac{1}{4}\,,\label{eq:check1}\\[4pt]
  \frac{\Delta}{\delta} &= 5\,.\label{eq:check2}
\end{align}
Both are exact rational numbers, involving no approximation.  The first
identity connects to the cone-partition geometry ($\sin^2 30° = 1/4$), while
the second connects to the golden ratio ($\sqrt{5}$ enters $\varphi$).  The
fact that a single pair of rationals encodes both the discrete icosahedral
symmetry and the continuous rotational geometry of $S^3$ is a non-trivial
consistency requirement.

Additionally, the numerator of $D^*$ satisfies
\begin{equation}
  67 = 24\!\times\! D^* = 24 \times \frac{67}{24}\,,
\end{equation}
which is a prime number.  The primality of $67$ ensures that $D^*$ does not
simplify to a ratio with a smaller denominator, thereby fixing the lattice
scale uniquely.

%======================================================================
\section{Dedekind Eta and the Bosonic String}\label{sec:dedekind}
%======================================================================

The Dedekind eta function is defined by
\begin{equation}\label{eq:eta}
  \eta(\tau) = q^{1/24}\prod_{n=1}^{\infty}(1-q^n)\,,
  \qquad q = e^{2\pi i\tau}\,,
\end{equation}
where $\tau$ is the modular parameter of the torus.  The exponent $1/24$
appearing in the leading $q$-expansion is precisely the elementary lattice
tension $\delta = 1/24$ of Axiom~II.  This is not a coincidence: the Dedekind
eta encodes the partition function of a single free boson on a torus, and the
$1/24$ arises from the Casimir energy
\begin{equation}\label{eq:casimir}
  E_0 = -\frac{1}{24}
\end{equation}
of each transverse oscillator mode.

The bosonic string requires $D_{\text{bos}} = 26$ spacetime dimensions for
Weyl anomaly cancellation.  In the SDGFT framework, this number emerges from
the relation
\begin{equation}\label{eq:dbos}
  D_{\text{bos}} = 2 + \frac{1}{\delta} = 2 + 24 = 26\,,
\end{equation}
where the additive~$2$ accounts for the longitudinal (light-cone) directions.
The 24~transverse dimensions of the bosonic string are thus in one-to-one
correspondence with the 24~cells (or, equivalently, the 24~vertices) of the
24-cell polytope.

The $\eta$-function furthermore enters the one-loop vacuum amplitude of the
bosonic string as $|\eta(\tau)|^{-48}$, where the exponent $48 = 2 \times 24$
equals the number of transverse oscillators (left + right movers).  In SDGFT,
$48 = 1/\delta + 1/\delta$.

The modular transformation $\tau \to -1/\tau$ acts on $\eta$ as
\begin{equation}\label{eq:modular}
  \eta(-1/\tau) = \sqrt{-i\tau}\;\eta(\tau)\,,
\end{equation}
which encodes T-duality in string theory and, within SDGFT, the self-duality of
the 24-cell: the polytope is invariant under the exchange of vertices and cells,
$f_0 \leftrightarrow f_3$, mirroring the modular $S$-transformation.

%======================================================================
\section{The Zero-Parameter Philosophy}\label{sec:zero}
%======================================================================

SDGFT occupies a unique position among candidate theories of quantum gravity by
containing \emph{zero} adjustable parameters.  Table~\ref{tab:compare}
contrasts this with the two main competitors.

\begin{table}[H]
\centering
\caption{Free parameters in candidate quantum-gravity
         frameworks.}\label{tab:compare}
\begin{tabular}{@{}lcc@{}}
\toprule
Framework & Free parameters & Landscape size \\
\midrule
String theory & moduli, fluxes & $\sim10^{500}$ \\
Loop quantum gravity & $\gamma_{\text{BI}}$ & 1 \\
SDGFT & 0 & 1 \\
\bottomrule
\end{tabular}
\end{table}

In string theory, the vacuum degeneracy arises because the six (or seven)
compact dimensions can be stabilised in exponentially many ways.  Each choice
yields a different low-energy physics, and no dynamical selection principle has
been identified~\cite{Bousso2000,Susskind2003}.

In LQG, the Barbero--Immirzi parameter $\gamma_{\text{BI}}$ is a one-parameter
ambiguity in the quantisation of the Ashtekar connection.  Its value is
traditionally fixed by demanding agreement with the Bekenstein--Hawking entropy
$S = A/(4G)$~\cite{Immirzi1997,Rovelli1996}, but this amounts to fitting one
number rather than predicting it.

SDGFT resolves both issues simultaneously.  The two axioms $\Delta$ and
$\delta$ are not fitted---they are \emph{derived} from the combinatorics of the
24-cell.  There is exactly one 24-cell, one value of $F_5$, and one minimal
lattice tension.  The absence of free parameters implies that every prediction
of the theory is a sharp numerical value, falsifiable by a single discrepant
measurement.

%======================================================================
\section{Conclusions}\label{sec:conclusions}
%======================================================================

We have laid the axiomatic foundations of Scale-Dependent Geometric Field
Theory.  The entire framework rests on two axioms:
\begin{enumerate}
  \item $\Delta = 5/24$: the Fibonacci-lattice conflict parameter, encoding the
    incompatibility of icosahedral and crystallographic symmetry on the 24-cell;
  \item $\delta = 1/24$: the elementary lattice tension, the minimal curvature
    quantum of the 24-cell.
\end{enumerate}
From these, the golden ratio $\varphi$, the cone half-angle $\theta_{\max} =
30°$, the six-cone partition, and the critical dimension $D^* = 67/24$ all
follow without free parameters.

The connection to the Dedekind eta function identifies $\delta$ with the
Casimir energy per bosonic oscillator and reproduces the bosonic-string
dimension $D_{\text{bos}} = 26$.  The self-duality of the 24-cell mirrors the
modular $S$-transformation.

Companion papers develop the $F_4$ symmetry and topology of the 24-cell
(Ref.~\cite{Besemer2026_02}) and the six-cone partition and flavour topology
(Ref.~\cite{Besemer2026_03}).

\begin{thebibliography}{99}

\bibitem{Bousso2000}
R.~Bousso and J.~Polchinski,
``Quantization of four-form fluxes and dynamical neutralization of the
cosmological constant,''
J.\ High Energy Phys.\ \textbf{06}, 006 (2000).

\bibitem{Susskind2003}
L.~Susskind,
``The anthropic landscape of string theory,''
arXiv:hep-th/0302219 (2003).

\bibitem{Immirzi1997}
G.~Immirzi,
``Quantum gravity and Regge calculus,''
Nucl.\ Phys.\ B (Proc.\ Suppl.) \textbf{57}, 65 (1997).

\bibitem{Rovelli1996}
C.~Rovelli,
``Black hole entropy from loop quantum gravity,''
Phys.\ Rev.\ Lett.\ \textbf{77}, 3288 (1996).

\bibitem{Coxeter1973}
H.~S.~M.~Coxeter,
\emph{Regular Polytopes}, 3rd ed.\ (Dover, New York, 1973).

\bibitem{Conway1988}
J.~H.~Conway and N.~J.~A.~Sloane,
``The cell structures of certain lattices,''
in \emph{Miscellanea Mathematica} (Springer, Berlin, 1988).

\bibitem{Besemer2026_02}
D.~A.~Besemer,
``Topology of the 24-Cell and $F_4$ Symmetry in Scale-Dependent Geometric
Field Theory,'' SDGFT-2026-02.

\bibitem{Besemer2026_03}
D.~A.~Besemer,
``Six-Cone Partition of the Three-Sphere and Flavour Topology,''
SDGFT-2026-03.

\end{thebibliography}

\end{document}
